\documentclass[11pt]{article}
\usepackage{amsmath,amssymb,amsthm}
\newtheorem{theorem}{Theorem}
\newcommand{\E}{\mathbb E}
\newcommand{\Cov}{\operatorname{Cov}}
\newcommand{\Var}{\operatorname{Var}}
\title{Depth-uniform stopped covariance and complete currents on one identified native fiber}\author{}\date{}
\begin{document}\maketitle
\paragraph{Identified input, not a generic effective measure.}
Use exactly the native probability of \cite[Section~3]{ND}:
\[
 dK=\widehat D^{-1}1_{\Omega_R}e^{-U(B)}dB,\qquad
 \Omega_R=\prod_{e=1}^q\{|B_e|_*<R\},\quad n=(\dim G)q,
\]
where $U=\frac12B^TQB-\ell\cdot B-W_{\rm nl}(B)-\log J(hB)$ is the actual
admitted density and $D^2U\ge\kappa_0 I>0$ on this original convex domain.
The criterion is the explicit one in (3.3), with the actual fixed-component coordinate count.
The sharp selector, genuine denominator, drift and Haar Jacobian are retained.
The supplement's convex-domain Poincare proof gives
\[
 \Var_KF\le\kappa_0^{-1}K|\nabla_BF|^2.
\]
We use this only on that identified component and on its genuine internal coordinate fibers.
It is not imposed on the arbitrary compact Wilson marginals of the inherited effective score calculation through the odd torus transfer moment calculation.

\begin{theorem}[Stopped means and all internal increments]
Let $F,H\in H^1(K)$ be actual carrying functions and let
$\mathcal F_0\supset\cdots\supset\mathcal F_m$ be any decreasing family of observations
on this same probability space, with $F_j=\E_K(F|\mathcal F_j)$ and $F_0=F$.
Then uniformly in $m$ and the observations,
\begin{align}
 \Var_KF_m&\le\kappa_0^{-1}K|\nabla F|^2,\\
 |\Cov_K(F_m,H_m)|&\le\kappa_0^{-1}
       (K|\nabla F|^2)^{1/2}(K|\nabla H|^2)^{1/2},\\
 \sum_{j=1}^m\left|K[(F_{j-1}-F_j)(H_{j-1}-H_j)]\right|
 &\le\kappa_0^{-1}(K|\nabla F|^2)^{1/2}(K|\nabla H|^2)^{1/2}.
\end{align}
Thus both the stopped inherited covariance and the accumulated scalar covariance increments
have a bound independent of the number of internal eliminations.
\end{theorem}
\begin{proof}
Conditional expectation is the orthogonal $L^2(K)$ projection. Total variance bounds
$\Var F_m$ by $\Var F$, then Poincare bounds the latter. Reverse-martingale increments
are orthogonal, so their summed squared norms are at most $\Var F$.
Cauchy--Schwarz over the increment index and Poincare give the third bound.
No spatial-overlap count or replacement of the conditional law has been used.
\end{proof}

\paragraph{Complete normalized auxiliary operations.}
The scalar theorem also survives adding a genuine source-free positive normalized kernel
$\Pi(d\xi|B)$ while keeping the same carrying function $F(B)$: its original variance is
still $\Var_KF$, and every actual output conditional mean is an $L^2$ contraction of it.
For a sequence of such complete Markov operations, conditional-variance differences telescope
and the same bound controls the sum of their covariance increments. This follows from the
Markov tower identity, even if successive output sigma-algebras are not literally nested.
It does not apply to a separately selected signed kernel, a source-zero tilt changing the
old marginal, or a changed carrier without its additional response. Current derivatives of
these auxiliary kernels still appear in the full weak transport; the scalar contraction does
not authorize deleting them before the complete normalization and integration.

\paragraph{Actual physical H and lower-row specialization.}
When the original complete top-scale interior and response hypotheses of supplement
Section5 hold, put $u(B)=\epsilon_{\rm bg}+h|B|_2$ (the subscript is a norm, not a square).
That section proves
\[
 |\nabla O_H(f)|\le C\|f\|_\infty\gamma^{-2}V_{D,f}h u(B),
\]
and for the literal ordinary partner, with its Casimir constant retained,
\[
 |\nabla r_H(f)|\le C\|f\|_\infty\gamma^{-2}
 [V_{D,f}+\beta_s^{-1}a^{-2}V_{Y,f}]h u(B).
\]
Define $A_f=\|f\|_\infty\gamma^{-2}V_{D,f}$ for $O_H$, or the same expression with
the bracket for $r_H$. Then every bound in the theorem is at most
\[
 C A_f A_k\,\kappa_0^{-1}h^2
       [\epsilon_{\rm bg}^2+h^2M_2],\qquad
 M_2=2n/\kappa_0+L_0^2/\kappa_0^2.
\]
This is the supplement's physical covariance scale now shown to control the entire
internal stopped-mean/increment decomposition of that same carrier. Saved couplings,
old mesh powers, fine support volumes and the actual target dimension are unchanged.
The absolute Casimir constant has not disappeared from raw masses merely because its
gradient is zero. No new claim of nondegeneracy or universal operator normalization follows.

\paragraph{Coordinate subfibers: conditional covariance and a bounded current.}
Split native coordinates as $B=(x,z)$ by complete original coordinate groups, so
$\Omega_R=\Omega_x\times\Omega_z$. The actual conditional density is
\[
 dK_x(z)=D(x)^{-1}1_{\Omega_z}e^{-U(x,z)}dz.
\]
The Hessian $D_z^2U\ge\kappa_0I$ and the ORIGINAL convex cutoff remain. Thus
\[
 \Var_{K_x}F\le\kappa_0^{-1}K_x|\nabla_zF|^2.
\]
Conditional Cauchy--Schwarz and Jensen imply, for every fixed $p\ge1$,
\[
 \|\Cov(F,H\mid x)\|_{L^p(\nu_x)}
 \le\kappa_0^{-1}\|\nabla_zF\|_{L^{2p}(K)}
                         \|\nabla_zH\|_{L^{2p}(K)}.
\]
The radial moment recursion (3.8) supplies the required higher gradient moments for the
actual physical H and lower rows above, retaining $n,L_0,\kappa_0$ explicitly.
No analogous pointwise conditional Poincare claim is made for a nonlinear observation.

For a fixed $x$, set $R_F=F-K_xF$ and solve the weak Neumann problem
\[
 -L_x\psi=R_F,\qquad L_x=\Delta_z-\nabla_zU\cdot\nabla_z,
 \qquad \partial_n\psi=0,\qquad K_x\psi=0.
\]
The conditional Poincare inequality gives existence and uniqueness in the mean-zero
Sobolev space. Put $W_F=\nabla_z\psi$, with zero components in $x$.
Then, as distributions INCLUDING the sharp boundary,
\[
 R_F\,1_{\Omega_z}e^{-U}
 =-\operatorname{div}_z(1_{\Omega_z}e^{-U}W_F).
\]
Its boundary partner vanishes by the constructed weak zero-normal-flux condition.
This does not assert zero boundary flux for an original physical current.
The energy estimate gives
\begin{align}
 K_x|W_F|^2&\le\kappa_0^{-1}\Var_{K_x}F,\\
 K|W_F|^2&\le\kappa_0^{-2}K|\nabla_zF|^2.
\end{align}
This replaces the exponential inverse-gap estimate of the compact AI41 construction on
this IDENTIFIED native sector. It is an $L^2$ current bound, not spatial locality, complete
current-jet control, or a bound on derivatives of $\psi$ with respect to $x$.

\paragraph{Whole source pair and exact compensation.}
Let $F_i=R_i+S_i$, with $S_i=K_xF_i$ and currents $W_i$ as above. The observation is
literally $x$, so $DqW_i=0$ pointwise. Conditional integration by parts gives
$K_x(W_iH)=K_x(R_iH)$ for admissible carrying functions.
For smooth carriers, the integrable degree-two scalar compensation in this alternative
weak source decomposition is
\[
 b_{12}=A_{12}+\tfrac12(W_1R_2+W_2R_1)+W_1S_2+W_2S_1.
\]
Here $W_iS_j=0$ because $S_j$ depends only on $x$, and
\[
 K_xb_{12}=K_xA_{12}+\Cov(F_1,F_2\mid x).
\]
These equalities are justified by the displayed weak adjoint pairings. Full WC11
transport against a varying tested kernel additionally needs the current derivatives and
carrying responses; their uniform bounds have not been inferred from this $L^2$ estimate.
At original XP3, $A_{12}=0$; the covariance is retained and quantitatively bounded,
not canceled by normalization. Subtracting it as an explicit degree-two source is a separate
prescription. Original independent scalar factors and raw denominators stay attached.

For an actual native readout $U_B$, testing the same distributional identity with
$\Phi(U_B)$ pushes the current forward as the vector-valued measure with coefficient
$D_BU_B\,W_i$. This is a valid FORWARD current transport with the actual map and measure;
it does not identify that current with an old Wilson EOM current or construct an inverse
physical-current lift. Constraints, chart Jacobians and any later kernel responses must
remain in their actual representation. No product of unrelated boundary distributions occurs.

\paragraph{The precise remaining identification gap.}
The result controls arbitrary stopped observations INSIDE ONE native component once its
actual carrier has the proved gradient norm. General RG histories change components,
references, selectors and carriers, and the old reconstructed word need not be the bare
physical H evaluated on this interior family. Our original-Wilson conditional mean
$\E_\mu(G|U_R)$ has gradient estimates under a DIFFERENT marginal law; these cannot
be inserted into $K$-Poincare without identifying its native pullback and controlling the
same gradient under $K$. A protected row independent of native coordinates has zero
variance exactly, but that does not extend to touched rows.
Thus this is a genuine depth-uniform scalar/current mechanism on an identified fiber,
not a solution of the full mixed-history MT36, the live-reference lift, or S94.
\begin{thebibliography}{9}
\bibitem{ND} \emph{Nonlinear dispersion and unequal spatial weights}, \texttt{lattice\_theory/block\_and\_transfer/conditional\_estimates/nonlinear\_dispersion\_and\_unequal\_spatial\_weights.tex} in this collection.
\end{thebibliography}
\end{document}
